\documentclass[11pt,a4paper]{article}
\usepackage[margin=23mm]{geometry}
\usepackage{amsmath,amssymb,booktabs,array}
\input{glyphtounicode}
\pdfgentounicode=1
\usepackage[hidelinks]{hyperref}
\setlength{\parindent}{0pt}
\setlength{\parskip}{6pt}
\begin{document}
\begin{center}
{\Large H1 and T1: a complete conditional lane-drop member}\\[4pt]
14 October 2026 --- consolidated for the original VSL-B2-G1 confirmation gate
\end{center}

The user selected T1's finite physical terminal interpretation, fixed
preparation handoff and restricted own-lane-first member. This note assembles
the exact upstream and terminal laws for contract confirmation. It closes
the bounded forward-model evidence review, not gradient or empirical
qualification. S13 remains unchanged.

\section*{1. Geometry, intervention, information and units}
A fixed road of length $L$ has preparation length $L_p=L-b$ and an initially
empty terminal length $b>0$. An initial uniform footprint $[0,\ell]$ lies
upstream, with $0<\ell\leq L_p$. There are no later external arrivals.
The three densities are $e_0,a_0,b_0\geq0$, with $a_0=b_0$ and
\[
 \rho_0=e_0+a_0+b_0,\qquad M=\ell\rho_0>0,\qquad r=e_0+a_0.
 \tag{1}
\]
$E$ is the ending lane, $A$ the adjacent continuing lane, and $B$ the other
continuing lane. The scalar $b_0$ is a density; $b$ is a length.

One constant common action $v\in[v_-,v_+]$, $0<v_-<v_+$, sets preparation
speed and the terminal free-flow parameter. Terminal queued motion need not
equal $v$. All geometry, initial traffic, $w,k_j,h,\gamma,\alpha,R,R_B$
are fixed across actions and observers. Use the declared triangular family
\[
 k_c(v)=\frac{wk_j}{v+w},\qquad C(v)=v k_c(v),\qquad w,k_j,h>0.
 \tag{2}
\]
The opportunity scale $h$ is physical and fixed, not a numerical time step.
Early and late work weights $\gamma,\alpha\geq1$ are distinct parameters.
$R,R_B>0$ are fixed downstream work ceilings, not future receiving traces.
Ordinary mass uses one normalized work unit.

\begin{center}
\begin{tabular}{lll}
\toprule
Quantity & Meaning & Units\\
\midrule
$L,L_p,b,\ell,h,z$ & lengths and physical progress & m\\
$v,w$ & free-flow and backward-wave speeds & m/s\\
$e,a,b_0,k_j,k_c$ & lane densities & veh-eq/m\\
$M,F_j,P_j,Y_j,N_j,q_j$ & mass or cumulative mass & veh-eq\\
$C,S,R,R_B$ & normalized receiving-work rates & work/s\\
$f,g_E,u_j$ & transfer-equivalent or exit mass rates & veh-eq/s\\
$m=f/h$ & lateral density source & veh-eq/(m s)\\
$\gamma,\alpha$ & work weights per transferred mass & dimensionless\\
$\mathrm{TST}$ & complete population-time integral & veh-eq s\\
\bottomrule
\end{tabular}
\end{center}

\section*{2. H1: state to spare service to accepted early transfer}
Initial densities satisfy $e_0,a_0,b_0\leq k_c(v_+)$.
Inside the remaining homogeneous occupied body,
\[
 S(v,a)=\min\{C(v),w(k_j-a)\},\quad
 f=\min\left\{ve,\frac{S(v,a)-va}{\gamma}\right\},\quad
 \dot e=-\frac f h,\quad \dot a=\frac f h,\quad \dot b_0=0.
 \tag{3}
\]
Own flow $va$ claims receiving first; $ve$ is the eager rate offer.
$f$ is an opportunity-equivalent bound for the same-coordinate lateral
source, not a second longitudinal boundary flow. Actual transfer over
$dx\,dt$ is $(f/h)dx\,dt$, not $\gamma$ times that mass.

In the invariant free-flow branch $e+a=r$ and $S=C$, so
\[
 \dot e=-\frac v h\min\left\{e,\frac{k_c(v)-r+e}{\gamma}\right\},
 \qquad a=r-e,\qquad b(t)=b_0.
 \tag{4}
\]
At $e=0$ transfer vanishes; at $a=k_c(v)$ spare service vanishes.
The locally Lipschitz scalar equation therefore has a unique solution with
\[
 0\leq e(t)\leq e_0,\qquad
 a_0\leq a(t)\leq\min\{r,k_c(v)\},\qquad b_0\leq k_c(v).
 \tag{5}
\]
Uniform initial characteristics experience the same local source until they
leave preparation. Consequently the remaining upstream body stays uniform
and its edges move at $v$, provided the terminal entrance accepts it.
No homogeneity is asserted inside the terminal.

With $z=vt$, $\delta(v)=k_c(v)-r$,
\[
 e_z=-\frac1h\min\{e,(\delta+e)/\gamma\}.
 \tag{6}
\]
The explicit speed factor cancels. Speed-sensitive preparation must come
through the declared normalized headroom, not elapsed time alone.
Donor-limited and service-limited branch references from $(z_s,e_s)$ are
\[
 e(z)=e_s e^{-(z-z_s)/h},\qquad
 e(z)=(e_s+\delta)e^{-(z-z_s)/(\gamma h)}-\delta,
 \tag{7}
\]
used only on their respective active branches.

\section*{3. The fixed handoff and conservative upstream balance}
Early lateral sources stop at $L_p$. No further lateral source operates
inside the terminal section; unfinished mass merges only at $L$.
The initial terminal histories are zero.
The preparation front and rear cross the handoff at
\[
 t_f^p=(L_p-\ell)/v,\qquad t_r^p=L_p/v.
\]
Terminal entry rates are, almost everywhere,
\[
 \dot F_E=ve(t),\quad \dot F_A=va(t),\quad \dot F_B=vb_0
 \quad\hbox{for }t_f^p<t<t_r^p,
 \tag{8}
\]
and zero otherwise. Single instants at interval boundaries do not change
the continuous-mass histories.

Let $s(t)$ be occupied length still in preparation: it equals $\ell$ before
$t_f^p$, $L_p-vt$ during the exit interval, and zero after $t_r^p$.
Stop the early source integral once $s=0$ and put
$J(t)=\int_0^t s(u)f(u)/h\,du$. Then
\[
 \begin{split}
 s\rho_0+F_E+F_A+F_B&=M,\\
 se+F_E+J&=\ell e_0,\\
 sa+F_A-J&=\ell a_0,\qquad sb_0+F_B=\ell b_0.
 \end{split}
 \tag{9}
\]
These identities follow by differentiating on each interval and checking
initial values. Already completed early transfers enter $F_A$ once;
they are not forced to merge again.

\section*{4. T1: physical storage, eligible demand and terminal receiving}
Define delays $\tau=b/v$, $\tau_b=b/w$, and per-lane jam storage $H=k_jb$.
Histories are extended as zero for negative times. For $j=E,A,B$,
\[
 P_j(t)=F_j(t-\tau),\quad
 N_j(t)=F_j(t)-Y_j(t),\quad
 q_j(t)=P_j(t)-Y_j(t).
 \tag{10}
\]
$P_j$ is cumulative minimum-travel eligibility, $N_j$ physical occupied
stock, and $q_j$ eligible but unserved mass. The eligible queue is not
the physical stock or the time when a traveler meets the queue's tail.
Boundary counts retain the terminal LWR travel/storage meaning without
independent interior position states.

Empty-link LTM receiving is at capacity while vacancy is positive:
\[
 V_j(t)=
 \begin{cases}
 k_jwt-F_j(t),&0<t\leq\tau_b,\\
 Y_j(t-\tau_b)+H-F_j(t),&t>\tau_b .
 \end{cases}
 \tag{11}
\]
Use the conservative all-action domain
\[
 M<H,\qquad 0\leq \dot F_j(t)\leq C(v).
 \tag{12}
\]
After $\tau_b$, $V_j\geq H-M>0$. Before it,
$F_j(t)\leq C(v)t<k_jwt$. The entrance therefore never holds H1.
This is a sufficient finite-storage proof, not an infinite buffer, a
queue-only storage bound, or a general spillback solver.

At the ending-lane exit, demand is zero before $\tau$.
Afterwards it is $\dot P_E$ when $q_E=0$ and $C(v)$ when $q_E>0$.
Both output histories obey minimum travel and finite sending.
The present node is not claimed to be Jin's general invariant junction model.

\section*{5. The selected restricted compulsory server}
The consolidated elementary member uses
\[
 v_+r<R\leq C(v_-),\qquad
 v_+b_0\leq R_B\leq C(v_-).
 \tag{13}
\]
The upper ceilings are a conservative qualification refinement of K2:
even the slowest free-flow lane can accommodate the declared exit budgets.
The lower bound for $R$ fits the entire ordinary $E+A$ load.
Any terminal queue is therefore caused by extra forced work, not hidden
ordinary overload.

Continuing traffic is served immediately:
$Y_A=P_A$, $Y_B=P_B$, $u_A=\dot P_A$.
Available ending service and actual flow are
\[
 c_E(t)=\min\left\{C(v),\frac{R-u_A(t)}{\alpha}\right\}>0,\qquad
 g_E=
 \begin{cases}
 c_E(t),&q_E>0,\\
 \min\{\dot P_E,c_E(t)\},&q_E=0 ,
 \end{cases}
 \quad \dot Y_E=g_E.
 \tag{14}
\]
Thus $\dot q_E=\dot P_E-g_E$, $q_E\geq0$, $Y_E\leq P_E$, and
\[
 u_A+\alpha g_E\leq R,\qquad
 u_A+g_E\leq R\leq C(v).
 \tag{15}
\]
Exactly $g_Edt$ exits, not $\alpha g_Edt$. There is no ending-lane forward
escape, hidden second queue, lost residual mass or terminal re-transfer.
On (13), the sending minimum in (14) is nonbinding because
$(R-u_A)/\alpha\leq R\leq C(v_-)\leq C(v)$.
It is retained to show the physical sending interface.

\textbf{Well-posedness of this member.}
The H1 solution fixes the absolutely continuous $F_j$ and $P_j$.
Let $p=P_E$, $K(t)=\int_0^t c_E(u)\,du$, and $X(t)=p(t)-K(t)$.
The unique reflected fluid-queue solution is
\[
 q_E(t)=X(t)-\min_{0\leq u\leq t}X(u),\qquad Y_E(t)=p(t)-q_E(t).
 \tag{16}
\]
$X(0)=0$; the running minimum prevents service of unavailable mass.
Differentiating almost everywhere gives (14). This proves existence and
uniqueness for the restricted boundary-coupled member, without invoking
a theorem for arbitrary weighted network nodes.

All mass is eligible by $t_*=L/v$. Thereafter $u_A=0$ and
$c_E=R/\alpha>0$, so
\[
 T_{\mathrm{bound}}=\frac Lv+\frac{\alpha q_E(L/v)}R.
 \tag{17}
\]
This equals full clearance when a positive continuing tail reaches $L/v$,
or when $q_E(L/v)>0$; otherwise it is an upper bound.
Using the bound in population-time integrals merely adds zero after
the actual last exit.

\section*{6. One joint domain with active early service and positive forced waiting}
The joint construction need not make early receiving irrelevant.
The following scaled algebraic family demonstrates compatibility; it
is not a calibrated scenario, selected traffic parameter set or rollout.
Take arbitrary reference speed $u>0$ and density $\kappa>0$, and set
\[
 \begin{gathered}
 v_-=u,\quad v_+=2u,\quad w=2u,\quad k_j=3\kappa,\\
 e_0=\frac7{10}\kappa,\quad a_0=b_0=\frac14\kappa,\quad
 \gamma=\alpha=2,\quad
 R=\frac{39}{20}\kappa u,\quad R_B=\kappa u,\\
 b=\ell>0,\quad L_p=\ell+h\log2,\quad L=L_p+b,\quad h>0.
 \end{gathered}
 \tag{18}
\]
Then $k_c(v_+)=3\kappa/2$, $C(v_-)=2\kappa u$,
$C(v_+)=3\kappa u$, $r=19\kappa/20$, and
\[
 M=\frac65\kappa\ell <3\kappa\ell=H,\qquad
 v_+r=\frac{19}{10}\kappa u<R<C(v_-),\qquad
 v_+b_0=\frac12\kappa u<R_B<C(v_-).
 \tag{19}
\]
Every initial density meets (5); invariant lane densities bound every entry
rate by $C(v)$ over the whole action interval.

At the faster action, the initial normalized accepted early rate is
\[
 \min\left\{\frac7{10},\frac{3/2-1/4}{2}\right\}\kappa
 =\frac58\kappa<\frac7{10}\kappa.
 \tag{20}
\]
Hence adjacent receiving limits the eager request. At the slower
action the service alternative is $7\kappa/8$, so donor availability binds.
The actions have different active preparation branches.

At the first terminal entry, exposure is $z_f=h\log2$.
The donor bound gives $e_f\geq e_0e^{-z_f/h}=7\kappa/20$.
For $v_+$, the initial service branch switches at
$e=11\kappa/20$ after $z_s=2h\log(25/22)<h\log2$.
Equation (7) gives the exact front residuals
\[
 e_f(v_-)=\frac7{20}\kappa,\qquad
 e_f(v_+)=\frac{11}{20}\kappa
     e^{-[h\log2-2h\log(25/22)]/h}
     =\frac{125}{352}\kappa>\frac7{20}\kappa.
 \tag{21}
\]
This is genuine, if modest, preparation dependence, not a speed-independent
progress law.

The first fast eligible pulse has weighted demand
\[
 v_+(a_f+\alpha e_f)=2u(r+e_f)
 \geq\frac{13}{5}\kappa u>R.
 \tag{22}
\]
Positive queue growth therefore persists over a nonzero interval by
continuity. Ordinary traffic still fits $R$, and (19) keeps terminal
receiving nonbinding at the upstream interface. Active early limitation,
positive forced waiting and protected homogeneous preparation coexist.
No full-TST preference, interior optimum or observer collision is inferred.

\section*{7. Complete conservation and the population-time objective}
For all times, with $s=0$ after preparation empties,
\[
 s(t)\rho_0+\sum_jN_j(t)+\sum_jY_j(t)=M .
 \tag{23}
\]
This follows from (9) and (10); work weights do not enter population counts.
All mass exits by (17). Complete TST and mean time are
\[
 \mathrm{TST}(v)=\int_0^{T_{\mathrm{bound}}}
       \left[M-\sum_jY_j(t)\right]dt,\qquad
 \overline T=\mathrm{TST}/M.
 \tag{24}
\]
Every initial position $x\in[0,\ell]$ has no-delay exit time $(L-x)/v$.
Only the ending stream has excess terminal delay. Since
$\sum_jP_j$ counts all no-delay eligibility,
\[
 M-\sum_jY_j=M-\sum_jP_j+q_E,\qquad
 \mathrm{TST}=M\frac{L-\ell/2}{v}+\int_0^{T_{\mathrm{bound}}}q_E(t)\,dt.
 \tag{25}
\]
This identity includes both sections' travel and every unit's excess waiting.
Eligible queue area measures excess time without equating $q_E$ with spatial
occupancy or moving the physical queue beyond the road.

\textbf{Deletion check.} For $\alpha=1$, (13) admits all combined ordinary
arrival flow. Thus $q_E=0$, $Y_j=P_j$, and (25) gives the common fastest optimum.
The same limit holds when the weighted arrivals already fit the node.
This is a declared limiting reduction, not general aggregate equivalence.

\textbf{Terminal-only reference.} Ordinary and ending masses four each,
uniformly eligible over two seconds, with $R=4$, $\alpha=2$ and slack
sending, give ending queue peak two, drain time one, queue area three and
complete-time sum eleven. The same masses eligible over four seconds give
zero queue and complete-time sum sixteen. These are isolated exact
references reused from K2, not H1 traffic experiments.

\section*{8. What information and speed sensitivity can be claimed}
Known $\ell$, initial symmetry and empty terminal histories give
\[
 e_0=E_0/\ell,\qquad a_0=b_0=(M-E_0)/(2\ell).
 \tag{26}
\]
Equations (4), (8), (10), (14) then determine the entire forward model.
Aggregate mass plus $E_0$ is sufficient. Queues and boundary histories are
predicted internal states, not additional initial measurements or future
service oracles. No discrete order is needed among this homogeneous
single-destination forced class.

For unchanged preparation labels per initial position, fixed $R,\alpha$
and the nonbinding-sending member, faster arrival cannot worsen complete
TST. Put $W=P_A+\alpha P_E$ and $Z=Y_A+\alpha Y_E$.
The work-conserving server gives
\[
 Z(t)=\inf_{0\leq u\leq t}\{W(u)+R(t-u)\},\qquad
 Y_A+Y_E+Y_B=Z/\alpha+(1-1/\alpha)P_A+P_B.
 \tag{27}
\]
Earlier cumulative arrival makes $W,P_A,P_B$ pointwise no smaller, hence
departures pointwise no smaller. In the full H1 model preparation may change
with advice, so this comparison does not settle its optimum.
An interior optimum is not guaranteed by queueing.

The equations are piecewise smooth with preparation branch switches,
action-dependent arrival endpoints, queue onset/emptying and capacity/domain
boundaries. This note qualifies neither differentiated event handling nor
autograd, finite differences, independently useful directions or a new
information/action certificate. Those are original Pass 3 obligations.

\section*{9. Assumption register, provenance and gate disposition}
\begin{itemize}
\item Retained task restrictions: finite cohort, initial continuing symmetry,
one common fixed advice, ideal admitted free-flow response, continuous mass,
fixed topology, full clearance and complete TST.
\item H1 adaptations: a uniform occupied footprint, a uniform lateral source
at the same road coordinate, an eager offer tied to physical progress,
scale $h$, the selected triangular speed family and unrestricted free-flow
admission to the terminal.
\item T1 selected meaning: finite terminal reach, initially empty histories,
no internal lateral sources after the handoff, point compulsory merging,
strict own-lane-first priority and one delayed class.
\item Declared terminal adaptations: constant $R,R_B,\alpha$, conservative
storage and ordinary-load restrictions, and all-action receiving ceilings
in (13). None is a field-calibrated zipper or synchronized-flow law.
\end{itemize}

Carey, Balijepalli and Watling (2013), \emph{Extending the Cell Transmission
Model to Multiple Lanes and Lane-Changing}, Section 5, supports local R1 allocation,
weighted receiving and a proportional rival, not (3)'s continuous source or
the complete compulsory server.
\href{https://eprints.whiterose.ac.uk/id/eprint/77315/7/Carey\%20et\%20al\%202013_with_coversheet.pdf}{Primary manuscript}.

Jin, \emph{Continuous formulations and analytical properties of the link
transmission model}, published 2015, primary preprint arXiv:1405.7080,
Sections 3.1--3.3, supports empty-link queue, vacancy, travel and storage
relations from boundary histories.
Its general invariant junction result is not attributed to the weighted
node here. \href{https://arxiv.org/pdf/1405.7080}{Primary preprint}.

Yperman, Logghe and Immers (2005), \emph{The Link Transmission Model: An
Efficient Implementation of the Kinematic Wave Theory in Traffic Networks},
Sections 3--4, supports link/node separation and the need for a stated merge
allocation. It does not calibrate this lane-end work factor.
\href{https://www.mech.kuleuven.be/cib/verkeer/dwn/pub/P2005C.pdf}{Primary paper}.

\textbf{Gate disposition.} The original Pass 2 evidence review is complete.
A unique, conservative, finite-clearance joined analytical member exists
with genuine early service limitation and forced waiting. The user selected
T1; this assembled exact contract awaits confirmation before freeze.
Applied MP6, observer matching, gradient correctness and independent
usefulness remain unestablished. If storage/admission, ordinary-load,
initial-symmetry or response conditions fail, reopen the smallest affected
G2/G3/G4 contract; do not clip mass, add queues or change laws silently.
No implementation, installation, traffic experiment or next review has begun.
\end{document}
